% NOMAD-Handbuch – Notfall-Wissensserver für Zuhause
% Bauen: make -C de/handbuch   (benötigt xelatex, polyglossia, tikz, siunitx, ImageMagick)
\documentclass[fontsize=13pt,a4paper,oneside]{scrbook}
\usepackage{fontspec}
\usepackage{polyglossia}
\setmainlanguage[spelling=new]{german}
\setmainfont{Liberation Sans}
\setmonofont{Liberation Mono}[Scale=0.95]
\usepackage[top=2.2cm,bottom=2.4cm,left=2.4cm,right=2.2cm]{geometry}
\usepackage{xcolor,tikz,graphicx,booktabs,longtable,array,tabularx,colortbl,enumitem,fancyhdr,amssymb,caption,float,ragged2e,siunitx}
\usepackage[hidelinks]{hyperref}

% Ohne Silbentrennung (Muster fehlen im System): linksbündig ist für Leser ab 60 ohnehin ruhiger.
\RaggedRight
\setlength{\parindent}{0pt}\setlength{\parskip}{0.7em}
\linespread{1.18}
\sisetup{mode=text,locale=DE,group-separator={.},output-decimal-marker={,}}

\definecolor{nomadbraun}{HTML}{3B4021}
\definecolor{nomadhell}{HTML}{F4EDD8}
\definecolor{warnrot}{HTML}{9B1C1C}
\definecolor{hinweisblau}{HTML}{1F4E79}
\definecolor{gruen}{HTML}{2F6B2F}

\addtokomafont{chapter}{\color{nomadbraun}\sffamily\bfseries}
\addtokomafont{section}{\color{nomadbraun}\sffamily\bfseries}
\addtokomafont{subsection}{\color{nomadbraun}\sffamily\bfseries}
\renewcommand*{\chapterformat}{\thechapter\enskip}

\pagestyle{fancy}\fancyhf{}
\fancyhead[L]{\small NOMAD}
\fancyhead[R]{\small\nouppercase{\leftmark}}
\fancyfoot[C]{\small Seite \thepage}
\renewcommand{\headrulewidth}{0.4pt}
\fancypagestyle{plain}{\fancyhf{}\fancyfoot[C]{\small Seite \thepage}\renewcommand{\headrulewidth}{0pt}}

% ---------- Kästen ----------
% \kasten{Farbe}{Titel}{Inhalt}: einfacher Rahmen aus TikZ, bricht nicht über Seiten.
\newcommand{\kasten}[3]{\par\noindent\begin{tikzpicture}
\node[draw=#1,fill=#1!7,line width=1.4pt,rounded corners=3pt,inner sep=9pt,text width=\dimexpr\linewidth-18pt-2.8pt\relax,align=flush left]{%
{\sffamily\bfseries\color{#1}#2}\par\vspace{2pt}\RaggedRight #3};
\end{tikzpicture}\par}
\newcommand{\hinweis}[1]{\kasten{hinweisblau}{Hinweis}{#1}}
\newcommand{\achtung}[1]{\kasten{warnrot}{Achtung}{#1}}
\newcommand{\merke}[1]{\kasten{gruen}{Merken}{#1}}
% Befehl zum Abtippen
\newcommand{\befehl}[1]{\par\noindent\begin{tikzpicture}
\node[draw=black!70,fill=black!5,line width=0.8pt,inner sep=7pt,text width=\dimexpr\linewidth-14pt-1.6pt\relax,align=flush left]{\ttfamily #1};
\end{tikzpicture}\par}
\newcommand{\taste}[1]{\fbox{\strut\,\textsf{#1}\,}}
\newcommand{\menue}[1]{\textbf{\textsf{#1}}}
\newcommand{\abhaken}{$\square$\enspace}
\newcommand{\zeile}[1][6cm]{\rule[-2pt]{#1}{0.4pt}}
\newcommand{\feld}[1][4cm]{\raisebox{-3pt}{\rule{0pt}{14pt}}\rule[-3pt]{#1}{0.4pt}}

% Abbildungen: verkleinerte JPG-Kopien aus bilder/web (siehe Makefile)
\newcommand{\bild}[3][0.92]{\begin{figure}[H]\centering
\includegraphics[width=#1\linewidth]{bilder/web/#2}
\caption{#3}\end{figure}}
\captionsetup{font={small,sf},labelfont=bf,labelsep=colon}

\title{\sffamily\bfseries NOMAD-Handbuch}
\subtitle{Ein Wissensserver, der ohne Internet funktioniert}
\author{Deutsche Fassung von Project NOMAD}
\date{Stand: Oktober 2026 \,\textperiodcentered\, Software-Version 1.35.2}

\begin{document}
\frontmatter
\begin{titlepage}
\centering
\vspace*{2.5cm}
{\sffamily\bfseries\fontsize{40}{46}\selectfont\color{nomadbraun} NOMAD\par}
\vspace{0.4cm}
{\sffamily\bfseries\fontsize{24}{28}\selectfont Handbuch\par}
\vspace{0.8cm}
{\Large Ein Wissensserver für Zuhause, der auch ohne Internet funktioniert\par}
\vspace{2.5cm}
\begin{tikzpicture}
\node[draw=nomadbraun,line width=1.6pt,rounded corners=4pt,inner sep=14pt,fill=nomadhell,text width=11cm,align=flush left]{%
\textbf{Dieses Heft gehört zu:}\par\vspace{6pt}
Rechner / USB-SSD: \zeile[6.2cm]\par\vspace{8pt}
Aufbewahrungsort: \zeile[6.6cm]\par\vspace{8pt}
Eingerichtet am: \zeile[3cm]\hfill von: \zeile[3.2cm]};
\end{tikzpicture}
\vfill
{\small Stand Oktober 2026, Software-Version 1.35.2.\par
Inoffizielle deutsche Fassung von Project NOMAD (Crosstalk Solutions), Apache-Lizenz 2.0.\par
Quellen dieses Heftes: \texttt{github.com/huppiflupp/project-nomad-de}, Ordner \texttt{de/handbuch}.\par}
\end{titlepage}

\chapter*{So benutzen Sie dieses Heft}
\addcontentsline{toc}{chapter}{So benutzen Sie dieses Heft}
Dieses Heft ist zum \textbf{Ausdrucken} gedacht. Legen Sie es zur USB-SSD oder zum Notfall-Laptop. Wenn es darauf ankommt, hat man meistens alles vergessen, deshalb steht hier jeder Schritt, auch der kleinste.

\textbf{Was Sie brauchen:} einen Rechner (am besten mit Windows, den Sie nicht verändern müssen), eine USB-Festplatte oder USB-SSD mit mindestens 512~GB, eine Internetverbindung \emph{zum Einrichten} und etwas Zeit. Danach läuft alles ohne Internet.

\textbf{Wie das Heft aufgebaut ist:}
\begin{itemize}[leftmargin=1.2em,itemsep=2pt]
\item Kapitel~1 erklärt, was NOMAD ist und \emph{ab wann} es ohne Internet funktioniert.
\item Kapitel~2 bis~6 führen Sie von der leeren USB-SSD bis zum bestandenen Offline-Test.
\item Kapitel~7 ist die \textbf{Notfall-Karte}, eine Seite zum Heraustrennen.
\item Kapitel~8 rechnet den Stromverbrauch durch (Akku, Solar, Kosten).
\item Kapitel~9 enthält Tabellen zum Ausfüllen von Hand: Softwarestand, Inhalte, Zugangsdaten.
\item Kapitel~10 beschreibt die Notfall-KI, einen Sprachassistenten ohne Internet.
\end{itemize}

\kasten{nomadbraun}{Zeichen in diesem Heft}{%
\abhaken Kästchen zum Abhaken.\par
\taste{Enter} Eine Taste auf der Tastatur.\par
\menue{Einstellungen} Ein Knopf oder Menüpunkt auf dem Bildschirm.\par
\befehl{so sieht ein Befehl aus}
Befehle tippen Sie genau so ab, mit Leerzeichen, und bestätigen mit \taste{Enter}. Groß- und Kleinschreibung zählt.}

\hinweis{Ein Teil der Inhalte (Wikipedia, medizinische Nachschlagewerke, Karten) liegt \textbf{auf Englisch} vor. Die Oberfläche von NOMAD und dieses Heft sind deutsch. Mehr dazu im Anhang „Was NOMAD nicht kann“.}

\tableofcontents
\mainmatter
\chapter{Was ist NOMAD, und ab wann geht es ohne Internet?}

\section{Kurz erklärt}
NOMAD ist ein kleiner Server, der Wissen bereithält: Nachschlagewerke, Wikipedia, Anleitungen, Karten und auf Wunsch einen Sprachassistenten (KI). Sie bedienen ihn im Browser, wie eine Internetseite. Der Unterschied: Die Seite liegt auf \emph{Ihrem} Rechner. Ist das Wissen einmal geladen, brauchen Sie kein Internet mehr.

\bild[0.8]{img-boot2}{Die Startseite von NOMAD heißt „Kommandozentrale“. Von hier aus erreichen Sie alle Funktionen.}

\section{Der wichtigste Satz}
\merke{NOMAD ist erst offline-fähig, wenn Sie die Inhalte \textbf{vorher mit Internet geladen} haben. Ein frisch eingerichtetes NOMAD ohne geladene Inhalte ist eine leere Bibliothek.}

\section{Was geht wann?}
Die Tabelle zeigt, welche Stufe Sie erreicht haben und was dann ohne Internet möglich ist.

\renewcommand{\arraystretch}{1.35}
\begin{center}
\small
\begin{tabularx}{\linewidth}{>{\RaggedRight\arraybackslash}p{3.4cm}>{\RaggedRight\arraybackslash}X>{\RaggedRight\arraybackslash}X}
\toprule
\textbf{Stufe} & \textbf{Das geht ohne Internet} & \textbf{Das braucht Internet} \\
\midrule
\textbf{A. Gestartet} (USB-SSD gebootet, NOMAD eingerichtet) &
Rechner starten, Kommandozentrale öffnen, deutsche Anleitungen (Docs) lesen. Die Einrichtung selbst braucht kein Internet. &
Alles Weitere: Inhalte, Karten, KI-Modelle, Aktualisierungen. \\
\addlinespace
\textbf{B. Inhalte geladen} (Kapitel~5) &
Wikipedia, Nachschlagewerke, Anleitungen durchsuchen und lesen (Kiwix). Dies ist die Stufe, die Sie für den Notfall wollen. &
Neue Inhalte, neue Fassungen, Updates von NOMAD. \\
\addlinespace
\textbf{C. Karten geladen} &
Kartenansicht der geladenen Länder (zum Beispiel Deutschland, ab \SI{0,1}{GB}). &
Weitere Länder laden. \\
\addlinespace
\textbf{D. KI eingerichtet} (Kapitel~10) &
Fragen an den lokalen Sprachassistenten, ohne dass etwas ins Internet geht. &
Das Sprachmodell einmal herunterladen. \\
\bottomrule
\end{tabularx}
\end{center}

\hinweis{Eselsbrücke: Alles, was Sie \textbf{herunterladen}, braucht Internet. Alles, was Sie \textbf{lesen oder suchen}, nicht.}

\section{Wie viel Platz brauchen die Inhalte?}
NOMAD selbst belegt etwa 5~GB. Der Rest hängt davon ab, was Sie laden. Die Zahlen stammen aus den Katalogen von NOMAD (Stand Oktober 2026); die Stufen \emph{Standard} und \emph{Umfassend} enthalten jeweils die darunter liegende Stufe mit. NOMAD zeigt Größen mit Faktor 1024 an, deshalb stehen dort Zahlen, die rund 2\,\% kleiner sind als hier.

\textbf{Deutsche Inhalte (empfohlen):}

\renewcommand{\arraystretch}{1.25}
\begin{center}
\small
\begin{tabular}{@{}lrrr@{}}
\toprule
\textbf{Themenbereich} & \textbf{Basis} & \textbf{Standard} & \textbf{Umfassend} \\
\midrule
Medizin (Deutsch) & \SI{0,16}{GB} & \SI{0,20}{GB} & \SI{0,69}{GB} \\
Nachschlagen und Lernen (Deutsch) & \SI{0,26}{GB} & \SI{6,3}{GB} & \SI{18,6}{GB} \\
Alltag und Technik (Deutsch) & \SI{0,03}{GB} & \SI{0,21}{GB} & \SI{3,8}{GB} \\
\midrule
\textbf{Summe der deutschen Bereiche} & \textbf{\SI{0,45}{GB}} & \textbf{\SI{6,7}{GB}} & \textbf{\SI{23,1}{GB}} \\
\bottomrule
\end{tabular}
\end{center}

Dazu kommt die \textbf{deutsche Wikipedia} nach Wunsch: von \SI{137}{MB} (die wichtigsten Artikel, kurz) über \SI{1,2}{GB} (Beliebte Artikel) und \SI{18}{GB} (alles ohne Bilder) bis \SI{51}{GB} (alles mit Bildern); mehr dazu in Kapitel~5.

\textbf{Englische Inhalte} gibt es zusätzlich, vor allem für Überleben, Vorsorge und Medizin-Nachschlagewerke (Kategorien mit dem Zusatz „Englisch“). Alle englischen Bereiche zusammen sind \SI{6,3}{GB} (Essential), \SI{35,7}{GB} (Standard) und \SI{89,1}{GB} (Comprehensive).

\section{Wie lange dauert das Laden?}
Die Dauer hängt von Ihrem Internetanschluss ab. Die Tabelle zeigt die reine Übertragungszeit bei voller Geschwindigkeit:

\begin{center}\small
\begin{tabular}{@{}lrrr@{}}
\toprule
\textbf{Menge} & \textbf{16 Mbit/s (DSL klein)} & \textbf{50 Mbit/s} & \textbf{250 Mbit/s} \\
\midrule
\SI{0,45}{GB} (Deutsch, Basis) & 4 min & 1 min & 15 s \\
\SI{6,7}{GB} (Deutsch, Standard) & 56 min & 18 min & 4 min \\
\SI{23}{GB} (Deutsch, Umfassend) & 3,2 Std. & 1,0 Std. & 12 min \\
\SI{89}{GB} (Englisch, alles) & 12,4 Std. & 4,0 Std. & 48 min \\
\bottomrule
\end{tabular}
\end{center}
In der Praxis dauert es eher 20 bis 50\,\% länger. Der Rechner darf währenddessen nicht in den Ruhezustand gehen (im NOMAD-System ist er abgeschaltet).

\section{Zwei Wege zum Ziel}
\begin{description}[leftmargin=1.2em,style=nextline]
\item[Weg 2 – USB-SSD (empfohlen)] Sie schreiben ein fertiges System auf eine USB-SSD und starten Ihren Windows-PC davon. Windows und Ihre Daten auf der Festplatte des PCs bleiben unberührt. Siehe Kapitel~2.
\item[Weg 1 – eigener Rechner mit Ubuntu] Sie installieren Ubuntu Linux auf einem Rechner, der danach nur noch NOMAD betreibt (zum Beispiel ein altes Notebook). Siehe Kapitel~3 und~4.
\end{description}

\chapter{Weg 2: Die USB-SSD vorbereiten (unter Windows)}

Dieses Kapitel macht aus einer leeren USB-SSD ein startfähiges NOMAD-System. Ihr Windows-PC wird dabei \textbf{nicht verändert}: Sie starten ihn später nur von der USB-SSD statt von der eingebauten Festplatte. Ziehen Sie die SSD ab und starten neu, ist alles wie vorher.

\section{Das brauchen Sie}
\begin{itemize}[leftmargin=1.5em,itemsep=2pt,label={\abhaken}]
\item Einen Windows-PC mit Internet (zum Herunterladen). Zum Starten des NOMAD-Systems geht auch ein anderer PC, der ab 2015 gebaut wurde, mit mindestens 4~GB Arbeitsspeicher (besser 8~GB).
\item Eine \textbf{USB-SSD oder USB-Festplatte mit mindestens 512~GB} und USB~3 (blauer oder beschrifteter Anschluss „SS“). Ein normaler USB-Stick ist zu klein und zu langsam.
\item Etwa \SI{6}{GB} freien Platz auf dem Windows-PC für die heruntergeladene Datei.
\item Das Programm \textbf{balenaEtcher} (kostenlos, \texttt{etcher.balena.io}). Alternativ geht Rufus (\texttt{rufus.ie}).
\end{itemize}

\achtung{Beim Schreiben wird die USB-SSD \textbf{vollständig gelöscht}. Prüfen Sie, dass nichts Wichtiges darauf liegt, und wählen Sie im Programm unbedingt die richtige Platte. Wählen Sie nie die Platte, auf der Windows läuft.}

\section{Schritt 1: Das System herunterladen}
\begin{enumerate}[leftmargin=1.5em,itemsep=3pt]
\item Öffnen Sie im Browser die Seite \texttt{github.com/huppiflupp/project-nomad-de/releases}.
\item Laden Sie die Datei \texttt{nomad.img.xz} aus der neuesten Version herunter (etwa \SI{5}{GB}). Entpacken Sie sie \textbf{nicht}; Etcher kann sie so lesen.
\item Notieren Sie die dort angegebene Prüfsumme (SHA-256): \feld[8.5cm]
\end{enumerate}

\subsection*{Prüfen, ob die Datei vollständig ist (empfohlen)}
Öffnen Sie die Eingabeaufforderung (Windows-Taste, „cmd“ tippen, \taste{Enter}) und geben Sie ein, mit dem Pfad Ihrer Datei:
\befehl{certutil -hashfile "C:\textbackslash Users\textbackslash Name\textbackslash Downloads\textbackslash nomad.img.xz" SHA256}
Die lange Zahl, die erscheint, muss mit der Prüfsumme von der Webseite übereinstimmen. Weicht sie ab, laden Sie die Datei neu herunter.

\section{Schritt 2: Auf die USB-SSD schreiben}
\begin{enumerate}[leftmargin=1.5em,itemsep=3pt]
\item Schließen Sie die USB-SSD an. Starten Sie \menue{balenaEtcher}.
\item Klicken Sie auf \menue{Aus Datei flashen} und wählen Sie \texttt{nomad.img.xz}.
\item Klicken Sie auf \menue{Ziel wählen} und wählen Sie Ihre USB-SSD. Kontrollieren Sie die Größe: Es muss die Größe Ihrer SSD sein (zum Beispiel 500~GB oder 1~TB).
\item Klicken Sie auf \menue{Flash!}. Bestätigen Sie die Windows-Rückfrage zur Administrator-Berechtigung.
\item Warten Sie, bis „Flash erfolgreich“ erscheint (je nach SSD 10 bis 40 Minuten). Etcher prüft das Ergebnis danach automatisch.
\end{enumerate}

\hinweis{Meldet Windows nach dem Schreiben „Sie müssen den Datenträger formatieren“, klicken Sie auf \menue{Abbrechen}. Das System auf der SSD ist ein Linux-System, das Windows nicht lesen kann. Das ist richtig so.}

Mit \emph{Rufus} geht es ebenso: Datei wählen, Laufwerk wählen, auf \menue{Start} klicken und, falls gefragt, den \textbf{DD-Modus} (Abbild-Modus) wählen.

\section{Schritt 3: Den PC von der USB-SSD starten}
\begin{enumerate}[leftmargin=1.5em,itemsep=3pt]
\item PC \textbf{herunterfahren} (nicht „Neu starten“: Windows nutzt sonst manchmal einen Schnellstart, der das Starten von der SSD erschwert).
\item USB-SSD anschließen, PC einschalten und \textbf{sofort mehrfach} die Taste für das Startmenü drücken. Welche Taste das ist, hängt vom Hersteller ab:
\end{enumerate}

\begin{center}\small
\begin{tabular}{@{}ll@{}}
\toprule
\textbf{Hersteller} & \textbf{Taste für das Startmenü} \\
\midrule
Acer, Dell, Lenovo, Samsung & \taste{F12} \\
Asus & \taste{Esc} oder \taste{F8} \\
HP & \taste{F9} (bei manchen erst \taste{Esc}) \\
MSI, Gigabyte & \taste{F11} oder \taste{F12} \\
\bottomrule
\end{tabular}
\end{center}
Diese Angaben sind Erfahrungswerte; steht beim Einschalten eine Zeile wie „Press F12 for Boot Menu“ auf dem Bildschirm, gilt diese. Ihre Taste tragen Sie hier ein: \taste{\feld[1.2cm]}

\begin{enumerate}[leftmargin=1.5em,itemsep=3pt,resume]
\item Wählen Sie im Startmenü Ihre USB-SSD (oft mit „USB“ im Namen) und bestätigen Sie mit \taste{Enter}.
\item Nach etwa einer Minute erscheint der Desktop (Abbildung~\ref{fig:boot}). Eine Anmeldung ist nicht nötig. Nach einer weiteren Minute öffnet sich der Browser mit der NOMAD-Kommandozentrale.
\end{enumerate}

\bild[0.7]{img-boot1}{Das System ist gestartet. Das Fenster mit der NOMAD-Seite öffnet sich von allein, sobald NOMAD bereit ist.}\label{fig:boot}

\hinweis{Beim ersten Start richtet das System NOMAD ein und vergrößert sich auf die ganze Platte. Das braucht etwa eine bis zwei Minuten länger. Schalten Sie in dieser Zeit nicht aus. Auch ohne Internet funktioniert das.}

\section{Wenn der PC nicht startet}
\begin{itemize}[leftmargin=1.5em,itemsep=3pt]
\item \textbf{Die SSD erscheint nicht im Startmenü.} Anderen USB-Anschluss versuchen (am besten blau, direkt am PC, nicht an einem Hub). Im Startmenü auf „UEFI“ achten.
\item \textbf{Meldung zu „Secure Boot“.} Das System ist mit aktivem Secure Boot getestet und sollte so starten. Zeigt Ihr PC trotzdem eine Sicherheitswarnung, hilft es, im BIOS Secure Boot auszuschalten (Menü „Security“ oder „Boot“). Notieren Sie sich, was Sie geändert haben.
\item \textbf{Schwarzer Bildschirm, Fehlermeldung oder Absturz.} Notieren Sie Hersteller und Modell des PCs im Feld „Hardware“ in Kapitel~9 und probieren Sie einen anderen PC. Sehr neue Grafikchips und manche WLAN-Chips können Probleme machen. Dafür gibt es kein allgemeines Mittel.
\item \textbf{BitLocker-Abfrage von Windows.} Wenn Windows beim nächsten normalen Start nach einem Wiederherstellungsschlüssel fragt, weil Sie im BIOS etwas geändert haben, brauchen Sie den Schlüssel aus Ihrem Microsoft-Konto (\texttt{aka.ms/myrecoverykey}). Notieren Sie sich vorher, ob BitLocker aktiv ist.
\end{itemize}

\section{Nach dem Start: Passwort ändern}
Benutzer und Passwort des Systems sind \textbf{bei jeder USB-SSD gleich}: Benutzer \texttt{nomad}, Passwort \texttt{nomad}. Das Passwort wird nur bei Änderungen am System gebraucht, aber jeder, der das Heft kennt, kennt es auch. Ändern Sie es, sobald das System steht:
\begin{enumerate}[leftmargin=1.5em,itemsep=3pt]
\item Öffnen Sie ein Terminal: gleichzeitig \taste{Strg} + \taste{Alt} + \taste{T}.
\item Tippen Sie \texttt{passwd} und \taste{Enter}.
\item Geben Sie das alte Passwort (\texttt{nomad}) ein, dann zweimal das neue. Beim Tippen erscheint nichts auf dem Bildschirm. Das ist normal.
\item Tragen Sie das neue Passwort im Feld „Zugangsdaten“ in Kapitel~9 ein.
\end{enumerate}

\chapter{Weg 1: Ubuntu auf einem eigenen Rechner installieren}

Dieses Kapitel ist für Sie, wenn Sie \textbf{keine USB-SSD} nutzen, sondern einen Rechner dauerhaft für NOMAD einrichten (zum Beispiel ein altes Notebook). Alles, was auf dem Rechner ist, wird \textbf{gelöscht}. Haben Sie die USB-SSD aus Kapitel~2 benutzt, überspringen Sie dieses und das nächste Kapitel und machen in Kapitel~5 weiter.

\hinweis{Die Bilder zeigen die Installation von Ubuntu 26.04 LTS in einer Testumgebung. Auf Ihrem Rechner sehen Sie dieselben Fenster, nur mit Ihren Namen und Plattengrößen.}

\section{Ubuntu-Stick erstellen}
Laden Sie die Datei „Ubuntu Desktop 26.04 LTS“ von \texttt{ubuntu.com/download/desktop} herunter (etwa 6~GB) und schreiben Sie sie mit balenaEtcher auf einen USB-Stick ab 8~GB, genau wie in Kapitel~2 beschrieben. Starten Sie den Rechner vom Stick (Startmenü, Tasten siehe Kapitel~2).

\section{Installation}
\bild[0.75]{02-sprache-de}{Sprache wählen: \menue{Deutsch}.}
\bild[0.75]{04-tastatur}{Tastaturbelegung: \menue{Deutsch}.}
\bild[0.75]{05-netzwerk}{Internetverbindung herstellen (Kabel oder WLAN). Ohne Internet geht die Installation nicht weiter.}
\bild[0.75]{07-art-der-installation}{Art der Installation: die normale Installation wählen.}
\bild[0.75]{10-festplatte}{Festplatte: „Festplatte löschen und Ubuntu installieren“. \textbf{Das löscht alles auf dieser Platte.}}
\bild[0.75]{12-konto-ausgefuellt}{Benutzerkonto: Name, Rechnername und ein Passwort, das Sie sich merken. Tragen Sie es in Kapitel~9 ein.}
\bild[0.75]{14-zusammenfassung}{Zusammenfassung: Prüfen und auf \menue{Installieren} klicken.}
\bild[0.75]{16-installation-fertig}{Wenn die Installation fertig ist, auf \menue{Jetzt neu starten} klicken. Den Stick danach abziehen, wenn der Rechner dazu auffordert.}
\bild[0.75]{17-anmeldung}{Mit dem Passwort aus dem Benutzerkonto anmelden.}

Nach der Anmeldung folgen einige Willkommens-Fenster. Sie können alle mit \menue{Weiter} oder \menue{Überspringen} beantworten. Zusätzliche Programme (Ubuntu~Pro) brauchen Sie nicht.

\merke{Prüfen Sie, dass der Rechner im Internet ist, bevor Sie weitermachen: Öffnen Sie den Browser (Firefox) und rufen Sie eine beliebige Seite auf.}

\chapter{Weg 1: NOMAD installieren}

\hinweis{Haben Sie die fertige USB-SSD aus Kapitel~2 benutzt, ist NOMAD schon eingerichtet. Springen Sie zu Kapitel~5.}

Auf dem frisch installierten Ubuntu richten Sie NOMAD mit einem einzigen Befehl ein. Der Rechner braucht dafür Internet; die Installation lädt etwa 4~GB.

\section{Terminal öffnen}
Drücken Sie gleichzeitig \taste{Strg} + \taste{Alt} + \taste{T}. Es öffnet sich ein schwarzes Fenster, das Terminal.
\bild[0.75]{25-terminal}{Das Terminal. Hier tippen Sie die Befehle ein.}

\section{Befehle eingeben}
Tippen Sie die folgenden Zeilen \textbf{genau so} ein (oder schreiben Sie sie in der Reihenfolge ab) und bestätigen Sie jede mit \taste{Enter}. Das Passwort, nach dem das System fragt, ist das Passwort Ihres Benutzerkontos; beim Tippen erscheint nichts, das ist normal.

\befehl{sudo apt-get update \&\& sudo apt-get install -y curl}
Der zweite Befehl ist lang und darf im Heft in zwei Zeilen stehen; tippen Sie ihn \textbf{ohne Zeilenwechsel} ab.
\befehl{curl -fsSL https://raw.githubusercontent.com/\allowbreak huppiflupp/\allowbreak project-nomad-de/\allowbreak refs/\allowbreak heads/\allowbreak main/\allowbreak install/\allowbreak install\_nomad.sh -o install\_nomad.sh}
\befehl{sudo bash install\_nomad.sh}

\bild[0.75]{29-installer-start}{Der Installer meldet sich auf Deutsch und fragt, ob er fortfahren soll.}

\begin{itemize}[leftmargin=1.5em,itemsep=3pt]
\item Auf die Frage „Möchten Sie wirklich fortfahren? (j/N)“ antworten Sie mit \taste{j} und \taste{Enter}.
\item Auf die Frage nach der Lizenzvereinbarung antworten Sie ebenfalls mit \taste{j} und \taste{Enter}.
\item Dann läuft die Installation selbst (etwa 5 bis 15 Minuten, je nach Internet). Meldungen bleiben auf dem Bildschirm stehen. Schließen Sie das Fenster nicht.
\end{itemize}

\bild[0.75]{32-fertig}{Am Ende steht: „Die Installation von Project NOMAD wurde erfolgreich abgeschlossen!“ und die Adresse der Oberfläche. Ein Hinweis zur fehlenden NVIDIA-Grafikkarte ist unbedenklich. Danach können Sie das Terminal schließen.}

\section{NOMAD im Browser öffnen}
Öffnen Sie den Browser (Firefox) und tippen Sie in die Adresszeile:
\befehl{http://localhost:8080}
\bild[0.75]{34-nomad-start}{Die Startseite „Kommandozentrale“. Wenn Sie sie sehen, ist NOMAD eingerichtet.}

\hinweis{Das Wort \texttt{localhost} heißt „dieser Rechner“. Von einem anderen Gerät im selben WLAN erreichen Sie NOMAD über die Adresse des Rechners, zum Beispiel \texttt{http://192.168.1.20:8080}. Die Adresse finden Sie, indem Sie im Terminal \texttt{hostname -I} eingeben. Tragen Sie sie in Kapitel~9 ein.}

\section{Ein-/Ausschalten und Neustart}
NOMAD startet bei jedem Einschalten des Rechners von selbst. Falls es einmal nicht erreichbar ist, finden Sie die Handgriffe auf der Notfall-Karte (Kapitel~7).

\chapter{Inhalte laden – dafür brauchen Sie Internet}

Dies ist der Schritt, der aus einem leeren NOMAD ein Notfall-Wissensserver macht. Planen Sie Zeit ein: Das Laden dauert je nach Menge und Anschluss Stunden. Das System darf währenddessen nicht ausgeschaltet werden.

\kasten{nomadbraun}{Vorher prüfen}{%
\abhaken Der Rechner ist mit dem Internet verbunden (Kabel ist am sichersten).\par
\abhaken Das Netzteil steckt (nicht nur Akku).\par
\abhaken Freier Platz reicht für die geplante Menge (Tabelle in Kapitel~1). Bei der USB-SSD mit 512~GB sind alle deutschen Bereiche in „Standard“ (\SI{6,7}{GB}) und die deutsche Wikipedia „Beliebte Artikel“ (\SI{1,2}{GB}) bequem möglich.}

\section{Weg A: Der Schnellstart-Assistent (empfohlen)}
Klicken Sie auf der Startseite auf \menue{Schnellstart}. Der Assistent hat vier Schritte: \menue{Apps}, \menue{Karten}, \menue{Inhalte}, \menue{Überprüfen}.
\bild[0.85]{35-schnellstart-1}{Schnellstart, Schritt „Apps“: Wählen Sie die Funktionen, die Sie möchten. Für den Notfall-Wissensserver wählen Sie die \emph{Wissensbibliothek} (Wikipedia, Nachschlagewerke).}
\bild[0.85]{38-inhalte}{Schritt „Inhalte“: oben Wikipedia in mehreren Größen, darunter die Themenbereiche. Oben zeigt eine Leiste, wie viel Speicher Ihre Auswahl belegt.}
\bild[0.85]{39-medizin-stufen}{Jeder Themenbereich hat drei Stufen: \emph{Essential} (das Wichtigste, klein), \emph{Standard} und \emph{Comprehensive}. Die Größen stehen dabei.}
\bild[0.85]{41-ueberpruefen}{Schritt „Überprüfen“: Zusammenfassung mit der Gesamtgröße. Prüfen Sie sie, bevor Sie bestätigen.}

\merke{Beginnen Sie mit \textbf{Basis} in den drei deutschen Bereichen (zusammen etwa \SI{0,5}{GB}) und der deutschen Wikipedia „Schnellreferenz“ (\SI{0,14}{GB}). Das geht in wenigen Minuten und gibt Ihnen sofort ein arbeitsfähiges System. Danach, am besten über Nacht: „Standard“ und die deutsche Wikipedia „Beliebte Artikel“. Für Überleben und Vorsorge gibt es bisher nur englische Inhalte (Kategorie „Überleben \& Vorsorge (Englisch)“, etwa \SI{2,3}{GB} in der kleinsten Stufe).}

\section{Weg B: Einzelne Inhalte später nachladen}
Unter \menue{Einstellungen} $\rightarrow$ \menue{Content-Manager} finden Sie dieselben Kataloge. Dort laden Sie einzelne Inhalte nach oder löschen sie, wenn der Platz knapp wird.

\section{Wikipedia: Welche Größe?}
Wikipedia gibt es auf \textbf{Deutsch und Englisch}, jeweils in mehreren Größen. Die Größen sind die Anzeige in NOMAD. Es lässt sich genau eine Wikipedia auswählen.

\begin{center}\small
\renewcommand{\arraystretch}{1.25}
\begin{tabularx}{\linewidth}{@{}>{\RaggedRight\arraybackslash}p{5.6cm}rX@{}}
\toprule
\textbf{Auswahl} & \textbf{Größe} & \textbf{Geeignet für} \\
\midrule
Deutsche Wikipedia – Schnellreferenz & \SI{0,14}{GB} & die wichtigsten Artikel, kurz; schnelle Antworten \\
Deutsche Wikipedia – Beliebte Artikel & \SI{1,2}{GB} & wichtige Artikel in voller Länge, ohne Bilder \\
Deutsche Wikipedia – Komplett (Kompakt) & \SI{4,6}{GB} & alle Artikel, nur die Einleitung \\
Deutsche Wikipedia – Komplett (ohne Bilder) & \SI{18,1}{GB} & alle Artikel ohne Bilder \\
Deutsche Wikipedia – Komplett (Vollständig) & \SI{51,2}{GB} & alles mit Bildern; nur für große Platten \\
\midrule
Englische Wikipedia – Schnellreferenz bis Vollständig & \SIrange{0,3}{121}{GB} & dieselben Größenstufen, auf Englisch \\
\bottomrule
\end{tabularx}
\end{center}

Neben der Wikipedia enthalten die deutschen Kategorien unter anderem die Medizin-Enzyklopädie der Wikipedia (WikiMed), Wikibooks, Wikiversity, das deutsche Wörterbuch Wiktionary, das Kinderlexikon Klexikon, das Koch-Wiki, die iFixit-Reparaturanleitungen und die Bibliothek des Projekts Gutenberg.

\section{Karten}
Karten laden Sie unter \menue{Einstellungen} $\rightarrow$ \menue{Karten-Manager}. Dort gibt es zwei Wege:
\begin{description}[leftmargin=1.2em,style=nextline]
\item[Nach Land oder Region (empfohlen)] Knopf \menue{Länder auswählen}: Land ankreuzen, mit dem Regler die \emph{Zoomstufe} einstellen, den Knopf \menue{Download starten} drücken. NOMAD schneidet nur die gewählten Länder aus einer Weltkarte heraus und zeigt vorher die Größe an.
\item[Weltkarte] Knopf \menue{Weltkarte herunterladen}: der ganze Planet, etwa \SI{129}{GB}. Nur für große Platten.
\end{description}
Die Zoomstufe bestimmt den Detailgrad und damit die Größe. Beispiel Deutschland (Anzeige in NOMAD):

\begin{center}\small
\renewcommand{\arraystretch}{1.25}
\begin{tabular}{@{}rll@{}}
\toprule
\textbf{Zoomstufe} & \textbf{Detailgrad} & \textbf{Größe} \\
\midrule
10 & bis Stadtebene & \SI{0,12}{GB} (121~MB) \\
12 & Straßen größerer Orte & \SI{0,68}{GB} \\
15 & einzelne Straßen & \SI{5,6}{GB} \\
\bottomrule
\end{tabular}
\end{center}
\bild[0.85]{map01-karten-manager}{Der Karten-Manager: „Länder auswählen“ ist der empfohlene Weg.}
\bild[0.85]{map02-laender-waehlen}{Land ankreuzen, Zoomstufe wählen; unten steht die zu erwartende Größe.}
\bild[0.85]{map03-karte-deutschland}{Die Karte von Deutschland (Zoomstufe 10), offline angezeigt.}

\hinweis{Die Länderliste und die Beschriftung der Karte sind englisch (zum Beispiel „Germany“). Unter „Kuratierte Kartenregionen“ stehen nur US-Gebiete; für Deutschland nehmen Sie „Länder auswählen“.}

\section{Fortschritt beobachten}
\bild[0.85]{43-downloads}{Auf der Seite „Downloads“ sehen Sie, was gerade geladen wird und wie weit es ist. Nach einem Neustart geht das Laden an der Stelle weiter, an der es aufgehört hat.}

\section{Wann ist es fertig?}
Alles ist geladen, wenn die Downloads-Seite keinen Eintrag mehr zeigt und in der \menue{Wissensbibliothek} (Kiwix) alle gewählten Inhalte aufgelistet sind. Machen Sie danach \textbf{sofort den Offline-Test} (Kapitel~6), solange das Internet noch erreichbar ist.

\hinweis{Notieren Sie in Kapitel~9, welche Inhalte Sie geladen haben und mit welchem Datum. So wissen Sie später, wie aktuell Ihr Wissen ist.}

\chapter{Der Offline-Test}

Erst dieser Test zeigt, ob Ihr System im Notfall wirklich ohne Internet funktioniert. Machen Sie ihn \textbf{einmal gründlich, solange das Internet noch erreichbar ist}, damit Sie fehlende Inhalte nachladen können. Wiederholen Sie ihn nach jeder größeren Änderung (Kapitel~9: Tabelle „Offline-Test“).

\section{Vorbereitung}
\begin{itemize}[leftmargin=1.5em,itemsep=3pt,label={\abhaken}]
\item Alle Downloads sind fertig (Kapitel~5, Seite „Downloads“).
\item Sie haben die Notfall-Karte (Kapitel~7) zur Hand.
\end{itemize}

\section{Internet trennen}
\begin{enumerate}[leftmargin=1.5em,itemsep=3pt]
\item \textbf{Netzwerkkabel ziehen.} Bei WLAN: oben rechts im Desktop auf das Netzwerksymbol klicken und WLAN ausschalten.
\item \textbf{Beweis:} Öffnen Sie das Terminal (\taste{Strg}+\taste{Alt}+\taste{T}) und tippen Sie:
\befehl{ping -c 2 8.8.8.8}
Es muss „100\,\% packet loss“ erscheinen (Abbildung~\ref{fig:offline}). Sonst sind Sie noch im Internet.
\end{enumerate}
\bild[0.8]{44-offline-test-terminal}{Offline-Beweis im Terminal (hier zusätzlich mit der Liste der geladenen Inhaltsdateien).}\label{fig:offline}

\section{Neu starten ohne Internet}
Starten Sie den Rechner \textbf{neu, ohne Netzwerk}. Das ist der eigentliche Test: Im Notfall startet er genau so. Warten Sie etwa zwei Minuten nach dem Desktop.

\section{Die Checkliste}
\renewcommand{\arraystretch}{1.5}
\begin{center}
\begin{tabularx}{\linewidth}{@{}c>{\RaggedRight\arraybackslash}X>{\RaggedRight\arraybackslash}p{3.2cm}@{}}
\toprule
 & \textbf{Prüfung} & \textbf{Ergebnis} \\
\midrule
\abhaken & Der Rechner startet ohne Netzwerkkabel bis zum Desktop. & \\
\abhaken & Im Browser \texttt{http://localhost:8080} zeigt die Kommandozentrale. & \\
\abhaken & Unter \menue{Docs} lassen sich die deutschen Anleitungen lesen. & \\
\abhaken & Die \menue{Wissensbibliothek} (Kiwix, \texttt{http://localhost:8090}) zeigt alle geladenen Bücher. & \\
\abhaken & Eine Suche nach einem Begriff liefert einen Artikel (zum Beispiel „Wasser“ oder „water“). & \\
\abhaken & Wikipedia öffnet einen Artikel samt Seitenlinks. & \\
\abhaken & (falls geladen) Karten zeigen den gewählten Ausschnitt. & \\
\abhaken & (falls eingerichtet) Die KI antwortet auf eine Frage. & \\
\bottomrule
\end{tabularx}
\end{center}

\bild[0.85]{45-kiwix-offline}{Die Wissensbibliothek (Kiwix) ohne Internet: Alle geladenen Bücher sind da.}
\bild[0.85]{47-wikipedia-offline}{Ein Wikipedia-Artikel ohne Internetverbindung.}

\hinweis{Zeigt NOMAD Bilder oder Texte in \emph{anderer} Sprache als Deutsch, ist das richtig: Die geladenen Bücher sind überwiegend englisch (Anhang).}

\section{Wenn etwas nicht geht}
\begin{itemize}[leftmargin=1.5em,itemsep=3pt]
\item \textbf{Seite nicht erreichbar:} noch eine Minute warten, dann Notfall-Karte (Kapitel~7), Punkt „NOMAD starten“.
\item \textbf{Ein Buch fehlt:} Internet wieder anschließen, Download im Content-Manager wiederholen, Test erneut machen.
\item \textbf{Zeigt nur eine weiße Seite:} Browser mit \taste{Strg}+\taste{F5} neu laden.
\end{itemize}

\chapter{Die Notfall-Karte}
\thispagestyle{plain}
Diese Seite ist einzeln gedacht: Drucken Sie sie ein zweites Mal und legen Sie sie sichtbar zum Rechner.
\small

\begin{tikzpicture}
\node[draw=warnrot,line width=2.2pt,rounded corners=4pt,inner sep=12pt,text width=\dimexpr\linewidth-24pt-4.4pt\relax,align=flush left]{%
{\sffamily\bfseries\Large\color{warnrot} NOTFALL-KARTE NOMAD}\par\vspace{6pt}
\textbf{1. Adresse:} Im Browser (Firefox) \texttt{http://localhost:8080} eintippen.\par
Von einem anderen Gerät im selben WLAN: \texttt{http://} \feld[3.8cm] \texttt{:8080}\par\vspace{6pt}
\textbf{2. Der Rechner fährt nicht hoch:} USB-SSD fest eingesteckt? Anderer blauer USB-Anschluss. Startmenü-Taste: \taste{\feld[1.1cm]}\par\vspace{6pt}
\textbf{3. Seite lädt nicht, obwohl der Desktop da ist:} 2 Minuten warten. Dann Terminal (\taste{Strg}+\taste{Alt}+\taste{T}) und:\par
\befehl{sudo bash /opt/project-nomad/start\_nomad.sh}\par
Passwort eingeben (Tipp unsichtbar), \taste{Enter}.\par\vspace{6pt}
\textbf{4. Alles prüfen:} \befehl{sudo docker ps}
Es sollten Zeilen mit \texttt{nomad\_} erscheinen.\par\vspace{6pt}
\textbf{5. Ausschalten:} Oben rechts auf das Ein/Aus-Symbol, \menue{Ausschalten}. Nie die USB-SSD abziehen, solange der Rechner läuft.\par\vspace{6pt}
\textbf{6. Strom sparen im Notfall:} Bildschirm dunkel, WLAN/Bluetooth aus, nur bei Bedarf einschalten (Kapitel~8).\par\vspace{6pt}
\textbf{7. Wissen suchen:} \menue{Wissensbibliothek} (\texttt{http://localhost:8090}), Suchfeld links oben. Texte sind überwiegend englisch.\par\vspace{6pt}
\textbf{8. Passwort:} \feld[6cm] \quad (Ort der Notiz: \feld[3.8cm])\par\vspace{6pt}
\textbf{9. Hilfe:} \feld[10cm]};
\end{tikzpicture}

\vspace{0.5em}
Software-Version 1.35.2, Stand des Heftes Oktober 2026.

\chapter{Energie: Wie lange läuft das System?}

Im Notfall sind Strom und Zeit knapp. Dieses Kapitel zeigt, wie Sie mit wenigen Zahlen ausrechnen, wie lange ein Akku, eine Powerstation oder eine Solarfläche das System versorgt. Rechnen Sie mit Ihren eigenen Messwerten; die Beispielwerte in diesem Kapitel sind \textbf{Annahmen} und mit einem Stern (*) gekennzeichnet.

\section{Die drei Formeln}
\kasten{nomadbraun}{Formeln}{%
\textbf{Energie} in Wattstunden (Wh) $=$ Leistung in Watt (W) $\times$ Zeit in Stunden (h)\par\vspace{4pt}
\textbf{Laufzeit} in Stunden $=$ Akkuinhalt in Wh $\times$ 0,85 $\div$ Leistung in W\par
\textbf{Der Faktor 0,85} (Wirkungsgrad des Wandlers) gilt, wenn Sie über einen Wechselrichter oder ein Netzteil gehen.\par\vspace{4pt}
\textbf{Kosten in €} $=$ Energie in kWh $\times$ Strompreis in €/kWh\par
1 kWh $=$ 1000 Wh.}

\section{Wie viel Strom braucht Ihr System?}
Die Leistung hängt vom Rechner ab. Richtwerte* (Mittel im Leerlauf mit Bildschirm aus):

\begin{center}\small
\renewcommand{\arraystretch}{1.25}
\begin{tabular}{@{}lrr@{}}
\toprule
\textbf{Gerät} & \textbf{Leerlauf*} & \textbf{unter Last*} \\
\midrule
Mini-PC / Einplatinenrechner & \SIrange{5}{15}{W} & \SIrange{15}{30}{W} \\
Notebook (Bildschirm an) & \SIrange{10}{20}{W} & \SIrange{30}{60}{W} \\
Desktop-PC (ohne Grafikkarte) & \SIrange{30}{60}{W} & \SIrange{80}{150}{W} \\
Desktop-PC mit Grafikkarte (KI) & \SIrange{50}{80}{W} & \SIrange{200}{400}{W} \\
USB-SSD (zusätzlich) & \SI{1}{W} & \SI{3}{W} \\
WLAN-Router (zusätzlich) & \SIrange{6}{12}{W} & \SIrange{6}{12}{W} \\
\bottomrule
\end{tabular}
\end{center}

\hinweis{Messen Sie selbst: Ein Strommessgerät für die Steckdose kostet etwa 10~€ und zeigt die Leistung direkt in Watt. Tragen Sie die Werte in die Tabelle am Ende dieses Kapitels ein.}

\section{Beispielrechnungen (mit Annahmen*)}
\textbf{Annahme*:} Notebook mit NOMAD, Bildschirm meist aus, im Mittel \SI{20}{W}.

\begin{description}[leftmargin=1.2em,style=nextline]
\item[Wie viel Energie am Tag?] \SI{20}{W} $\times$ 24~h $=$ \SI{480}{Wh} $=$ \SI{0,48}{kWh}.
\item[Kosten bei Dauerbetrieb?] \SI{0,48}{kWh} $\times$ 0,35~€/kWh* $=$ 0,17~€ am Tag, rund 61~€ im Jahr.
\item[Laufzeit an einer Powerbank?] Eine Powerbank mit \SI{99}{Wh} (größte erlaubte Akkus im Flugzeug): $99 \times 0{,}85 \div 20 \approx$ \SI{4,2}{h}.
\item[Laufzeit an einer Powerstation?] Eine Powerstation mit \SI{500}{Wh}: $500 \times 0{,}85 \div 20 \approx$ \SI{21}{h}.
\item[Wie viel Solarfläche?] Ein Panel mit \SI{100}{W} liefert im Sommer rund \SI{350}{Wh} am Tag (5 volle Sonnenstunden*, 30\,\% Verluste: $100 \times 5 \times 0{,}7$), im Winter in Deutschland oft nur \SI{50}{Wh} bis \SI{100}{Wh} am Tag.
\end{description}

\begin{center}\small
\renewcommand{\arraystretch}{1.25}
\begin{tabular}{@{}lrrr@{}}
\toprule
\textbf{Pro Tag} & \textbf{20~W*, 24~h} & \textbf{20~W*, 4~h} & \textbf{60~W*, 4~h} \\
\midrule
Energiebedarf & \SI{480}{Wh} & \SI{80}{Wh} & \SI{240}{Wh} \\
Powerbank \SI{99}{Wh} reicht & 0,2 Tage & 1,1 Tage & 0,35 Tage \\
Powerstation \SI{500}{Wh} reicht & 0,9 Tage & 5,3 Tage & 1,8 Tage \\
Panel \SI{100}{W} deckt im Sommer* (\SI{350}{Wh}) & 73\,\% & 100\,\% & 100\,\% \\
Panel \SI{100}{W} deckt im Winter* (\SI{75}{Wh}) & 16\,\% & 94\,\% & 31\,\% \\
\bottomrule
\end{tabular}
\end{center}
\small „Reicht“ heißt: Akkuinhalt $\times$ 0,85 $\div$ Energiebedarf pro Tag. Die Laufzeit-Beispiele der Zeilen darüber gelten bei Dauerbetrieb.\normalsize

\section{Tipps zum Strom sparen}
\begin{itemize}[leftmargin=1.5em,itemsep=3pt]
\item \textbf{Bildschirm aus oder dunkel.} Bei Notebooks sind das oft 3 bis 6~W.
\item \textbf{Nur nutzen, wenn nötig.} Der Rechner lässt sich in Sekunden ausschalten; NOMAD startet beim Einschalten selbst.
\item \textbf{Keine KI nebenher.} Die KI (Kapitel~10) verbraucht je nach Rechner das Drei- bis Zehnfache.
\item \textbf{Router aus,} wenn Sie ihn nicht brauchen: Offline funktioniert alles auch am Rechner selbst.
\end{itemize}

\section{Ihre Messwerte}
\renewcommand{\arraystretch}{1.9}
\begin{center}
\begin{tabularx}{\linewidth}{@{}>{\RaggedRight\arraybackslash}p{4cm}XXX@{}}
\toprule
\textbf{Gerät / Zustand} & \textbf{Leistung (W)} & \textbf{Stunden/Tag} & \textbf{Wh/Tag} \\
\midrule
\feld[3.6cm] & & & \\
\midrule
\feld[3.6cm] & & & \\
\midrule
\feld[3.6cm] & & & \\
\midrule
\feld[3.6cm] & & & \\
\midrule
\multicolumn{3}{@{}r}{\textbf{Summe Wh/Tag:}} & \\
\bottomrule
\end{tabularx}
\end{center}
Akku/Powerstation: \feld[2cm]~Wh $\times$ 0,85 $\div$ Summe \feld[2cm]~W $=$ \feld[2cm]~Stunden.

\chapter{Tabellen zum Ausfüllen}

Füllen Sie diese Seiten von Hand aus, solange Sie das System vor sich haben. Später ist das die einzige Stelle, an der steht, was Sie eingerichtet haben.
\renewcommand{\arraystretch}{1.9}

\section{Hardware}
\begin{tabularx}{\linewidth}{@{}>{\RaggedRight\arraybackslash}p{4.8cm}X@{}}
\toprule
Rechner (Hersteller, Modell) & \\
\midrule
Arbeitsspeicher (GB) & \\
\midrule
USB-SSD (Hersteller, Größe) & \\
\midrule
Startmenü-Taste & \\
\midrule
Secure Boot im BIOS geändert? (ja/nein) & \\
\midrule
Netzteil / Ladegerät (Watt) & \\
\midrule
Wo liegt das Zubehör? & \\
\bottomrule
\end{tabularx}

\section{Softwarestand}
\begin{tabularx}{\linewidth}{@{}>{\RaggedRight\arraybackslash}p{4.4cm}XX@{}}
\toprule
 & \textbf{Beim Einrichten} & \textbf{Letzte Prüfung} \\
\midrule
Datum & & \\
\midrule
NOMAD-Version (Fußzeile der Startseite) & & \\
\midrule
Ubuntu-Version & & \\
\midrule
Image-Datei / Prüfsumme & & \\
\midrule
Letztes Update eingespielt am & & \\
\bottomrule
\end{tabularx}

\hinweis{Wie ein Update geht: Startseite $\rightarrow$ \menue{Einstellungen} $\rightarrow$ \menue{System-Update} $\rightarrow$ „Erneut prüfen“, bei Bedarf „Update starten“. Dafür braucht der Rechner Internet. Danach der Offline-Test (Kapitel~6).}

\section{Zugangsdaten}
\achtung{Dieses Heft ist ein Papier. Schreiben Sie nur Passwörter hinein, wenn es an einem sicheren Ort liegt.}
\begin{tabularx}{\linewidth}{@{}>{\RaggedRight\arraybackslash}p{4.8cm}X@{}}
\toprule
Benutzername des Systems & \\
\midrule
Passwort des Systems & \\
\midrule
WLAN-Name & \\
\midrule
WLAN-Passwort & \\
\midrule
Router-Adresse / Passwort & \\
\midrule
Adresse von NOMAD im Heimnetz & \\
\bottomrule
\end{tabularx}

\section{Geladene Inhalte}
\begin{tabularx}{\linewidth}{@{}>{\RaggedRight\arraybackslash}p{4.4cm}>{\RaggedRight\arraybackslash}p{2.2cm}X@{}}
\toprule
\textbf{Inhalt} & \textbf{Größe} & \textbf{Geladen am / Anmerkung} \\
\midrule
Wikipedia (Auswahl) & & \\
\midrule
Medizin & & \\
\midrule
Überleben und Vorsorge & & \\
\midrule
Bildung und Nachschlagen & & \\
\midrule
Heimwerken und Reparatur & & \\
\midrule
Landwirtschaft und Essen & & \\
\midrule
Computer und Technik & & \\
\midrule
Karten & & \\
\midrule
KI-Modell & & \\
\bottomrule
\end{tabularx}

\section{Sicherung und Offline-Test}
\begin{tabularx}{\linewidth}{@{}>{\RaggedRight\arraybackslash}p{3.2cm}XXX@{}}
\toprule
\textbf{Datum} & \textbf{Offline-Test bestanden?} & \textbf{Sicherung (Wo?)} & \textbf{Unterschrift} \\
\midrule
 & & & \\
\midrule
 & & & \\
\midrule
 & & & \\
\midrule
 & & & \\
\bottomrule
\end{tabularx}

\hinweis{Eine Sicherung (Backup) schützt gegen einen Ausfall der USB-SSD. Am einfachsten: Eine zweite USB-SSD mit demselben System (Kapitel~2) herstellen und die Inhalte nachladen. Das Heft hilft Ihnen dabei, den Stand nachzuvollziehen.}

\chapter{Die Notfall-KI: Ein Sprachassistent ohne Internet}

NOMAD kann einen Sprachassistenten (KI) betreiben, der \textbf{auf Ihrem Rechner} läuft: Sie stellen Fragen in ganzen Sätzen, die Antwort wird lokal berechnet. Es geht nichts ins Internet.

\kasten{warnrot}{Was die KI nicht ist}{%
Die KI kann sich irren und Dinge erfinden. Bei Medizin, Sicherheit, Recht und Chemikalien \textbf{prüfen Sie jede Antwort an den Quellen} in der Wissensbibliothek (Kapitel~5) oder in gedruckten Unterlagen. Sie ist ein Helfer zum Suchen und Erklären, kein Arzt.}

\section{Was dafür nötig ist}
\begin{itemize}[leftmargin=1.5em,itemsep=3pt]
\item \textbf{Internet beim Einrichten:} Das Sprachmodell (je nach Größe 1 bis mehrere GB) wird einmal heruntergeladen.
\item \textbf{Arbeitsspeicher:} Je größer das Modell, desto besser die Antworten, aber desto mehr Speicher und Zeit. Richtwerte*:
\end{itemize}

\begin{center}\small
\renewcommand{\arraystretch}{1.3}
\begin{tabularx}{\linewidth}{@{}>{\RaggedRight\arraybackslash}p{4.2cm}>{\RaggedRight\arraybackslash}p{3.5cm}X@{}}
\toprule
\textbf{Rechner} & \textbf{Modellgröße*} & \textbf{Zu erwarten*} \\
\midrule
Wenig Speicher (4--8~GB), keine Grafikkarte & ca. 1--3 Milliarden Parameter & einfache Fragen; langsam bis ordentlich \\
\midrule
Mittel (16~GB) & ca. 7--9 Milliarden & brauchbare deutsche Antworten, bei Rechnern ohne Grafikkarte einige Sekunden bis Minuten pro Antwort \\
\midrule
Grafikkarte (NVIDIA, 8~GB+ Speicher) & ca. 8--14 Milliarden & flott, bessere Qualität \\
\bottomrule
\end{tabularx}
\end{center}
\small * Diese Richtwerte stammen aus allgemeiner Erfahrung und sind von uns für dieses Heft noch \textbf{nicht} auf Ihrem Rechner gemessen. NOMAD empfiehlt selbst, mit kleinen Modellen zu beginnen und größere auszuprobieren.\normalsize

\section{Einrichten}
\begin{enumerate}[leftmargin=1.5em,itemsep=3pt]
\item Internet anschließen.
\item Startseite $\rightarrow$ \menue{Schnellstart} oder \menue{Einstellungen} $\rightarrow$ \menue{Apps verwalten}: den \textbf{KI-Assistenten} installieren.
\item In den KI-Einstellungen ein \textbf{kleines} Modell wählen und laden. Erst probieren, dann ein größeres.
\item Internet trennen, Frage stellen (Kapitel~6, Checkliste).
\end{enumerate}

\hinweis{Hat der Rechner eine NVIDIA-Grafikkarte, nutzt NOMAD sie, wenn der Installer sie erkennt (steht am Ende der Installation). Ohne Grafikkarte läuft die KI auf dem Prozessor und ist langsamer.}

\section{Eigene Dokumente für die KI}
In der Wissensdatenbank können Sie Dateien (PDF, Text) hochladen, aus denen die KI Antworten ableitet. Sinnvoll sind Ihre eigenen Notfallpläne, Telefonlisten, das Handbuch als PDF oder Medikamentenpläne der Familie.

\kasten{hinweisblau}{Noch nicht geprüft}{%
Ob die KI auch die per Kiwix geladenen Bücher (Wikipedia, Nachschlagewerke) \emph{selbst} nutzt, haben wir noch nicht abschließend geprüft. Verlassen Sie sich nicht darauf: Suchen Sie bei wichtigen Fragen selbst in der Wissensbibliothek.}

\section{Fragen stellen}
\begin{itemize}[leftmargin=1.5em,itemsep=3pt]
\item Kurz und konkret fragen: „Wie bereite ich Wasser mit Tabletten auf?“
\item Bei Zahlen und Dosierungen immer gegenprüfen.
\item Antwort unklar? Frage anders formulieren oder die Quelle in der Wissensbibliothek lesen.
\end{itemize}

\section{Ihre Notizen zur KI}
\renewcommand{\arraystretch}{1.9}
\begin{tabularx}{\linewidth}{@{}>{\RaggedRight\arraybackslash}p{4.8cm}X@{}}
\toprule
Gewähltes Modell & \\
\midrule
Größe auf der Platte & \\
\midrule
Dauer einer typischen Antwort & \\
\midrule
Hochgeladene Dokumente & \\
\bottomrule
\end{tabularx}

\appendix
\chapter{Was NOMAD nicht kann}

Damit Sie sich im Notfall nicht auf etwas verlassen, das nicht da ist, hier die bekannten Grenzen dieser deutschen Fassung (Stand Oktober 2026).

\begin{description}[leftmargin=1.2em,style=nextline]
\item[Manche Inhalte sind nur auf Englisch da.] Oberfläche, Installationsprogramm, Anleitungen und dieses Heft sind deutsch. Es gibt deutsche Wikipedia, Medizin-Enzyklopädie (WikiMed), Lehrbücher, Wörterbuch, Koch- und Reparaturanleitungen. Für \emph{Überleben und Vorsorge} und die medizinischen Nachschlagewerke (CDC, NHS, Militärmedizin) gibt es bisher nur englische Bücher. Deshalb sind die englischen Kategorien mit dem Zusatz „Englisch“ gekennzeichnet.
\item[Karten: Ortsnamen und Länderliste englisch.] Karten für Deutschland und jedes andere Land sind möglich (Kapitel~5). Die Auswahlliste der Länder und die Beschriftung der Karte („Germany“, „Cologne“) sind aber englisch. Die „kuratierten Kartenregionen“ im Karten-Manager enthalten nur Gebiete der USA; für Deutschland nehmen Sie „Länder auswählen“.
\item[Kein Windows-Programm.] NOMAD läuft unter Linux. Auf Windows ohne Linux gibt es nur eine Anleitung des Originals (WSL2) in englischer Sprache. Dieses Heft empfiehlt deshalb die USB-SSD.
\item[Kein Mac, keine ARM-Rechner.] Das System ist für gängige PCs (x86\_64) gebaut. Macs mit Apple-Chip und kleine ARM-Platinen (zum Beispiel Raspberry Pi) werden von dieser Version nicht unterstützt.
\item[Kein Passwortzwang.] Das System der USB-SSD hat bei jeder SSD dasselbe Anfangspasswort. Ändern Sie es (Kapitel~2).
\item[Fremde Hardware.] Ob das System auf Ihrem PC startet (Grafik, WLAN, Startmenü), hängt vom PC ab und lässt sich nicht für alle Geräte vorab zusichern.
\item[Updates.] NOMAD kann sich selbst aktualisieren, aber nur mit Internet; Inhalte veralten (Datum in Kapitel~9 notieren).
\end{description}

\end{document}
